\section{Introduction}

Large language models are increasingly used as judges for open-ended tasks.
Most work evaluates these judges at the final decision layer: whether the model
assigns the correct preference, score, or reward. This view misses a preceding
source of error. A judge can produce a plausible final decision while relying
on an incomplete set of evaluation dimensions. For example, a criteria list may check
surface helpfulness and factuality, yet omit instruction constraints,
evidence-grounding requirements, safety boundaries, or domain-specific failure
modes that human evaluators explicitly care about.

This is a different object of study from judge accuracy, preference alignment,
or reward-model ranking. Those metrics ask whether the final scalar or pairwise
decision agrees with a label. We ask whether the evaluative basis available to
the judge already contains the human-prioritized dimensions needed to make that
decision. If the basis is incomplete, a correct-looking decision can still hide
systematic missing dimensions.

We argue that this dimension-selection step should be studied as a first-class
research problem. The relevant question is not only ``is the judge accurate?'',
but also ``what dimensions does the judge consider before judging?'' We call
systematically omitted human evaluation dimensions \emph{evaluation blind
spots}. This framing moves LLM-as-a-Judge research upstream from score
calibration to the construction of the evaluative basis itself.

Open-ended criteria elicitation is difficult to optimize with standard supervised
or reinforcement learning recipes. Unlike math, code, or closed-form QA, there
is no unique target string. A good elicited criteria set may differ lexically
from the human-gold criteria while still covering the same semantic dimension; another
set may be longer but mostly redundant; a third may invent criteria that are
not grounded in the task. The training signal must therefore measure semantic
coverage, penalize duplicate dimensions, and reject unverifiable criteria.

We introduce Blind-Spot Coverage (BSC), a computable metric that measures the
fraction of human-gold evaluation dimensions semantically covered by generated
criteria. BSC is designed not only as an evaluation metric but also as a
verifiable reward. It rewards coverage of human dimensions, subtracts
redundancy among generated criteria, and penalizes invalid or hallucinated
criteria through a verifier. This converts open-ended criteria elicitation into a
reward-optimizable task without requiring an exact reference output.

Our framework, BlindSpot-RL, formulates a test of whether this reward can
extend RLVR/GRPO beyond deterministic domains into open-ended
criteria elicitation. The
training protocol separates human hard-gold evaluation data from proxy-gold
training data: RubricBench \texttt{test\_main} remains a held-out BSC benchmark,
while large-scale proxy criteria are elicited from non-overlapping query pools
with multi-teacher criteria elicitation and fail-closed verification.

This gives the paper a problem--metric--method structure rather than a claim
that a larger policy is simply better. Our contributions are:
\begin{itemize}
  \item We identify evaluation blind spots as an upstream failure mode of
  LLM-as-a-Judge: models may score plausibly while omitting human-prioritized
  dimensions.
  \item We propose Blind-Spot Coverage, a semantic metric and verifiable reward
  that combines gold-dimension coverage, redundancy control, and invalidity or
  hallucination filtering.
  \item We formulate an evidence-gated RLVR/GRPO training route for
  open-ended, semantic, subjective criteria elicitation by optimizing
  BSC rather than exact answers; empirical dimension-recovery wording is
  allowed only after hard-gold, downstream, ablation, and dimension-transition
  evidence gates pass.
\end{itemize}

The experiments are deliberately claim-gated: hard-gold holdout evaluation,
proxy-gold training isolation, downstream holdouts, confidence intervals,
ablations, human audit, semantic visualization, and submission-readiness gates
determine which empirical statements are safe to include.
